%% TEMPO: fleet-pooled learning for recourse in vehicle routing with
%% simultaneous pickup and delivery and stochastic demands.
%% Build: pdflatex tempo; bibtex tempo; pdflatex tempo; pdflatex tempo
%% Numbers marked \TBD are filled from scripts/tempo_pool_tables.py.
\newif\ifanon
\anonfalse

\documentclass[preprint,authoryear,12pt]{elsarticle}
\usepackage{amsmath,amssymb,amsthm}
\usepackage{setspace}
\usepackage{booktabs}
\usepackage{graphicx}
\usepackage{xcolor}
\usepackage[colorlinks=true,allcolors=blue!50!black]{hyperref}
\onehalfspacing

\newtheorem{assumption}{Assumption}
\newtheorem{lemma}{Lemma}
\newtheorem{theorem}{Theorem}
\newtheorem{proposition}{Proposition}
\newtheorem{remark}{Remark}
\newcommand{\TBD}{\textcolor{red}{[TBD]}}
\newcommand{\E}{\mathbb{E}}

\journal{Computers \& Operations Research}

\begin{document}
\begin{frontmatter}

\title{Learning the day as a fleet: pooled demand evidence for
recourse in vehicle routing with simultaneous pickup and delivery}

\ifanon
  \author{Anonymised for review}
\else
  \author{Quang-Vinh Dang}
  \ead{vinh.dq4@buv.edu.vn}
  \address{British University Vietnam, Hung Yen, Vietnam}
\fi

\begin{abstract}
Pickup volumes on a delivery day often share a common cause, such as
weather, a promotion or a local event. A vehicle that sees heavy pickups
at its first stops learns something about its later stops, and so do
the other vehicles of the fleet. Existing recourse policies for vehicle
routing with stochastic demands either assume the demand distribution is
known or learn it from the vehicle's own observations. We propose TEMPO
(Team Evidence and Multilevel Pooling for Online recourse), which pools
the pickups observed by the whole fleet in a two-level Bayesian model of
city-wide and zone-level day factors, and lets every vehicle choose
between continuing, handing its remaining route to a standby vehicle and
returning to the depot with the value tables of a policy library indexed
by the factor. We show that the excess cost of such a plug-in rule over
the library policy at the true factor is bounded by the squared error of
the factor estimate at the decision stops, and we give the point at which
pooling across zones starts to hurt. On \TBD{} instances from three
benchmark families, pooling closes \TBD{} of the gap between the policy
fitted under the planning law and the policy that knows the day's
factors, against \TBD{} for per-vehicle learning. The two-level model is
never worse than the better of full and zone-level pooling, and it keeps
its advantage when the factors drift during the day, when deliveries
carry their own factor, and when the pickup noise is misspecified.
\end{abstract}

\begin{keyword}
Vehicle routing \sep Stochastic demands \sep Recourse \sep Bayesian
learning \sep Information pooling
\end{keyword}
\end{frontmatter}

% ====================================================================
\section{Introduction}\label{sec:intro}

A vehicle that collects and delivers goods along a fixed route does not
know the pickup volumes in advance. Its load rises and falls along the
route, and a breach of capacity forces an expensive emergency vehicle.
Recourse policies decide at each stop whether to continue, to return to
the depot, or to hand the rest of the route to another vehicle
\citep{yang2000restocking,secomandi2009reopt,florio2022partial}. Almost
all of them assume that the probability law of the demands is known.

On a real day the law is not known. Demand volumes depend on the
weather, on promotions, on local events and on the day of the week, and
these causes act on many customers at once. A vehicle that meets heavy
pickups at its first three stops has learned that the day is heavy. The
other vehicles of the fleet have learned the same thing, if the cause is
shared. \citet{florio2023bayesian} model such a common factor for a
single route and update it by Bayes' rule as the vehicle observes its
customers. They report savings of more than ten per cent over policies
that assume independent demands. Their vehicle, however, learns only
from its own customers.

This paper asks what a fleet gains when its vehicles learn together. The
answer depends on how the day's variation is shared. If one city-wide
factor drives every customer, all vehicles should pool their evidence.
If each district has its own factor, pooling across districts mixes
unrelated evidence and can be worse than not pooling at all. Real days
lie in between. We therefore model the day by a city-wide factor and
zone-level deviations, estimate their variances from past days, and
compute the exact posterior of each zone's factor from the pickups the
fleet has seen so far. Each vehicle then uses the factor of the zones it
has still to visit.

The contributions are as follows.
\begin{enumerate}
\item A recourse method for the vehicle routing problem with
  simultaneous pickup and delivery and stochastic demands that pools
  demand evidence across the vehicles of a fleet. Decisions use a
  library of recourse policies indexed by the day factor, with exact
  value tables, so the learning layer adds no fitting during the day.
\item A regret bound: the excess cost of the plug-in rule over the
  library policy at the true factor is at most a sum, over decision
  stops, of local constants times the squared error of the factor
  estimate. The bound turns a statistical question (how well is the
  factor estimated at each stop?) into a cost statement, and it explains
  why pooling helps most at the first stops, where most recourse
  decisions are taken.
\item A threshold at which full pooling becomes worse than pooling
  within zones, in terms of the pickup noise, the zone-level variance
  and the number of zones.
\item A computational study on \TBD{} instances from the
  \citet{dethloff2001vehicle} and \citet{salhi1999cluster} benchmarks and
  on instances built on the road networks of five cities, with
  independent factor draws per instance, instance-clustered confidence
  intervals, and four forms of misspecification.
\end{enumerate}

The method sits on top of an execution policy for the same problem
\citep{dang2026baton}, which prices handoffs and depot returns from route
geometry and fits per-stop continuation costs by backward induction. Our
contribution is the learning layer, not the execution policy.

% ====================================================================
\section{Related work}\label{sec:related}

\paragraph{Recourse for stochastic demands} The vehicle routing problem
with stochastic demands assigns customers to a priori routes and
specifies what happens when a vehicle cannot serve a customer
\citep{bertsimas1992vrpsd,gendreau1996stochastic,oyola2018svrp1}.
\citet{yang2000restocking} add preventive returns to the depot, chosen by
dynamic programming over the vehicle's load; exact methods for the
resulting problem include \citet{louveaux2018exact} and the rule-based
and hybrid recourse policies of \citet{salavati2019rulebased,salavati2019hybrid}.
\citet{secomandi2001rollout} and \citet{secomandi2009reopt} re-optimise
the remaining route as demands are revealed, \citet{goodson2013rollout}
and \citet{goodson2016restocking} extend rollout policies to several
vehicles and duration limits, and \citet{florio2022partial} study partial
re-optimisation. All of these assume a known demand law.

\paragraph{Learning the demand law} \citet{florio2023bayesian} introduce
a common external factor that scales every customer's demand and update
its posterior along a fixed route. \citet{popovic1995bayesian} treats
uncertain demand in a Bayesian way at the planning stage. Our model
differs in two respects: the factor has a city-wide and a zone-level
part, and the evidence comes from the whole fleet.

\paragraph{Cooperation between vehicles} Paired recourse strategies let
two vehicles share customers when one runs out of capacity
\citep{ak2007paired,zhu2014paired}. They share work, not information.
Distributionally robust formulations protect a plan against an
ambiguous law instead of learning it \citep{ghosal2020dro}.

\paragraph{Pooling and shrinkage} Estimating many related means by
shrinking each towards a common value is the empirical Bayes idea of
\citet{efron1973stein}; hierarchical models choose the amount of
shrinkage from the data \citep{gelman2006prior}. Our regret bound uses
the margin argument of plug-in classification
\citep{audibert2007fast} and the performance-difference identity of
reinforcement learning \citep{kakade2002approx}.

% ====================================================================
\section{Problem and execution model}\label{sec:model}

\paragraph{Fleet and routes} A plan assigns the customers to $K$ routes
served from one depot. Route $r$ visits stops $k=1,\dots,m_r$ in a fixed
order. At stop $k$ the vehicle delivers $d_k$ and collects $p_k$; its
load changes by $g_k = p_k - d_k$. The departure load is the planned
delivery volume, and $B_r$ denotes the residual capacity at departure.
If the running load $W_k$ exceeds $B_r$, the vehicle breaches at stop
$k$: an emergency vehicle serves the rest of the route at price
$E_{r,k}$. After serving stop $k$ the vehicle may instead hand the rest
of the route to a standby vehicle at price $H_{r,k}$, or return to the
depot at price $R_{r,k}$, which resets the load to zero. Prices are built
from the route geometry and a three-class fleet cost model
\citep{dang2026baton}: planned vehicles, standby vehicles paid per use,
and emergency vehicles hired at surge prices.

\paragraph{Day factors} Customers lie in $Z$ geographic zones. On a
given day, zone $z$ has a pickup factor $a_z = \exp(b_z)$ with
\begin{equation}\label{eq:factor}
b_z = \mu + \eta_z, \qquad \mu \sim N(0, s_c^2), \qquad
\eta_z \sim N(0, s_z^2),
\end{equation}
independent. Given the factors, the pickup of customer $i$ in zone $z$
is gamma distributed with mean $a_z \bar p_i$ and coefficient of
variation $c$; deliveries follow the planning law. The planning law is
the special case $a_z = 1$ for all $z$.

\paragraph{Policy library} For each value $a$ on a grid $\mathcal{A}$,
$\pi_a$ is the recourse policy fitted by backward induction under
pickups scaled by $a$ \citep{dang2026baton}. For each $\pi_a$ we compute
exactly, on a load grid, the expected remaining cost $Q_a(s,x)$ of
taking action $x \in \{\text{continue},\text{handoff},\text{restock}\}$
in state $s=(r,k,W)$ and following $\pi_a$ afterwards. The library is
built once per plan, before the day.

% ====================================================================
\section{TEMPO: pooled learning during the day}\label{sec:method}

\paragraph{Posterior of the zone factors} Each pickup observed so far
contributes the gamma log-likelihood of $b_{z(i)}$. Under
\eqref{eq:factor}, the posterior of $b_z$ given all pickups observed by
the fleet is computed exactly on a grid of $b$ values: integrate the
city factor $\mu$ out of the other zones' likelihoods, convolve with the
zone-level prior, and multiply by zone $z$'s own likelihood. The cost is
a few one-dimensional convolutions per stop.

\paragraph{Variances from past days} The variances $s_c^2$ and $s_z^2$
are estimated by the method of moments from complete past days: the
spread of the zones' log ratios of observed to planned pickups within a
day estimates $s_z^2$ plus sampling noise, and the variance of the
daily average estimates $s_c^2$ plus a known share of the rest.

\paragraph{Decision rule} At stop $k$ of route $r$ the vehicle computes
its factor $\hat a_{r,k}$ as the planned-pickup-weighted mean of the
posterior means of the zones of its remaining customers, and takes
\begin{equation}\label{eq:rule}
x_{r,k} = \arg\min_x Q_{\hat a_{r,k}}(s_{r,k}, x).
\end{equation}
Special cases of the same rule are the policy fitted under the planning
law ($\hat a = 1$), per-vehicle learning (only the vehicle's own
pickups, as in \citealp{florio2023bayesian}), full pooling (one factor
for the fleet) and pooling within zones.

% ====================================================================
\section{Analysis}\label{sec:theory}

Fix a route and its true factor $a$, and write $e_k = \hat a_k - a$ for
the estimation error at decision stop $k$.

\begin{assumption}[Lipschitz in the factor]\label{ass:lip}
For every state $s$ and action $x$, $|Q_a(s,x) - Q_{a'}(s,x)| \le
\ell(s)\,|a-a'|$.
\end{assumption}

\begin{assumption}[Margin]\label{ass:margin}
Let $\Delta(s)$ be the difference between the second smallest and the
smallest of $Q_a(s,\cdot)$. For all $t \ge 0$,
$P(\Delta(s_k) \le t \mid e_k) \le f(s_k)\,t$ almost surely.
\end{assumption}

\begin{lemma}\label{lem:one}
Under Assumptions~\ref{ass:lip} and~\ref{ass:margin},
$\E[Q_a(s_k,x_k) - \min_x Q_a(s_k,x) \mid e_k] \le 4\,\ell(s_k)^2
f(s_k)\, e_k^2$.
\end{lemma}

\begin{proof}
If $x_k$ is not a minimiser of $Q_a(s_k,\cdot)$, then for a minimiser
$x^*$, $Q_a(s_k,x_k) \le Q_{\hat a}(s_k,x_k) + \ell|e_k| \le
Q_{\hat a}(s_k,x^*) + \ell|e_k| \le Q_a(s_k,x^*) + 2\ell|e_k|$. The loss
is therefore at most $2\ell|e_k|$, and a wrong choice requires
$\Delta(s_k) \le 2\ell|e_k|$. Taking the conditional expectation given
$e_k$ and applying Assumption~\ref{ass:margin} gives the bound.
\end{proof}

\begin{theorem}\label{thm:regret}
Let $J(\cdot)$ denote the expected cost of the route under the true
law, and let $L$ be the rule~\eqref{eq:rule}. Then
\[
J(L) - J(\pi_a) \le \sum_k \E\bigl[4\,\ell(s_k)^2 f(s_k)\, e_k^2\bigr],
\]
where the sum runs over the decision stops reached under $L$.
\end{theorem}

\begin{proof}
By the performance-difference identity \citep{kakade2002approx} for a
finite-horizon problem with benchmark $\pi_a$,
$J(L) - J(\pi_a) = \E_L \sum_k [Q_a(s_k,x_k) - Q_a(s_k,\pi_a(s_k))]$.
Each term equals $[Q_a(s_k,x_k) - \min_x Q_a(s_k,x)] + [\min_x
Q_a(s_k,x) - Q_a(s_k,\pi_a(s_k))]$; the second bracket is not positive,
and Lemma~\ref{lem:one} bounds the first.
\end{proof}

\begin{remark}
The benchmark is the library policy at the true factor. The reference
used in Section~\ref{sec:study}, which acts greedily with respect to
$Q_a$ at every stop, is a rollout of $\pi_a$ and is at least as good
\citep{bertsekas1997rollout}; Theorem~\ref{thm:regret} does not bound
the gap to it.
\end{remark}

\paragraph{Estimation error of the pooling rules} With $\kappa = 1/c^2$,
one pickup carries Fisher information $\kappa$ about $b_{z}$, whatever
its planned volume. With $n_z$ pickups observed in zone $z$, the zone's
evidence acts, to first order, as one observation of $b_z$ with variance
$c^2/n_z$. Writing $\Sigma = s_c^2\mathbf{1}\mathbf{1}^\top + s_z^2 I$
and $\Lambda = \mathrm{diag}(n_z/c^2)$, the two-level posterior has
error $[(\Sigma^{-1}+\Lambda)^{-1}]_{zz}$, the smallest of any rule that
uses the same data. Pooling within the zone has error
$(1/s^2 + n_z/c^2)^{-1}$ with $s^2 = s_c^2 + s_z^2$. Full pooling mixes
the other zones' deviations into the estimate of $b_z$.

\begin{proposition}\label{prop:threshold}
Suppose the $Z$ zones are observed equally, $n_z = n/Z$. As $n_z$ grows,
the error of full pooling tends to $s_z^2(1 - 1/Z) > 0$, while that of
pooling within the zone is of order $c^2/n_z$. To first order, pooling
within the zone is the better rule once
$n_z > c^2 / \bigl(s_z^2 (1 - 1/Z)\bigr)$.
\end{proposition}

With $c = 0.3$, $s_z = 0.2$ and two zones the threshold is 4.5 pickups
per zone. Most recourse decisions are taken at the first two or three
stops \citep{dang2026baton}, when few pickups have been seen; this is
where pooling matters and where full pooling can still beat pooling
within zones.

% ====================================================================
\section{Computational study}\label{sec:study}
\TBD

% ====================================================================
\section{Conclusion}\label{sec:conclusion}
\TBD

\bibliographystyle{elsarticle-harv}
\bibliography{tempo}
\end{document}
